% Part: many-valued-logic
% Chapter: sequent-calculus
% Section: introduction

\documentclass[../../../include/open-logic-section]{subfiles}

\begin{document}

\olfileid[tr]{mvl}{seq}{int}

\olsection{Giriş}

Klasik mantığın sekant hesabı etkili ve basit bir !!{derivation} sistemidir.
Bir çok değerli mantık sonlu sayıda doğruluk değerine sahip bir matrisle, yani
$V$ sonlu olacak biçimde tanımlanmışsa bu mantık için de bir sekant hesabı
verilebilir. Bunun nasıl yapılacağına ilişkin düşünce, klasik durumdaki
sekantların anlamları ile çıkarım kurallarının biçiminden gelir.

Şu sekantı hatırlayın:
\begin{align*}
    !A_1, \dots, !A_m & \Sequent !B_1, \dots, !B_n.
\intertext{Bu sekant şu !!{formula} olarak yorumlanabilir:}
    (!A_1 \land \cdots \land !A_m) & \lif (!B_1 \lor \cdots \lor
!B_n).
\end{align*}
Başka bir deyişle !!^a{valuation} $\pAssign{v}$, bir
$\Gamma\Sequent\Delta$ sekantını ancak ve ancak bazı $!A\in\Gamma$ için
$\pValue{v}(!A)=\False$ ya da bazı $!A\in\Delta$ için
$\pValue{v}(!A)=\True$ ise \emph{sağlar}. Bu yorumda başlangıç sekantları
$!A\Sequent !A$ her zaman sağlanır; çünkü ya $\pValue v(!A)=\True$ ya da
$\pValue v(!A)=\False$ olur.

Yan formüller $\Gamma,\Delta$ atıldığında $\Log{LK}$ içindeki koşul
bağlacının çıkarım kuralları şöyledir:

\begin{defish}
    \Axiom$ \fCenter !A$
    \Axiom$ !B \fCenter $
    \RightLabel{\LeftR{\lif}}
    \BinaryInf$ !A \lif !B \fCenter $
    \DisplayProof
    \hfill
    \Axiom$ !A \fCenter !B$
    \RightLabel{\RightR{\lif}}
    \UnaryInf$ \fCenter !A \lif !B $
    \DisplayProof
\end{defish}

Yukarıdaki semantik sekant yorumunu uygularsak $\LeftR{\lif}$ kuralı,
$\pValue v(!A)=\True$ ve $\pValue v(!B)=\False$ ise
$\pValue v(!A\lif !B)=\False$ olduğunu söyler. Benzer biçimde
$\RightR{\lif}$ kuralı, $\pValue v(!A)=\False$ ya da
$\pValue v(!B)=\True$ ise $\pValue v(!A\lif !B)=\True$ olduğunu söyler.
Aslında bu koşullar iki yönlüdür. Standart biçimlerindeki $\LeftR{\land}$ ve
$\RightR{\lor}$ kurallarına karşılık gelen koşullar iki yönlü olmazdı; ancak
bu kuralların iki yönlülüğü veren alternatif biçimleri vardır:

\begin{defish}
    \Axiom$!A, !B, \Gamma \fCenter \Delta$
    \RightLabel{\LeftR{\land}}
    \UnaryInf$!A \land !B, \Gamma \fCenter \Delta$
    \DisplayProof
    \hfill
    \Axiom$ \Gamma \fCenter \Delta, !A, !B$
    \RightLabel{\RightR{\lor}}
    \UnaryInf$ \Gamma \fCenter \Delta, !A \lor !B$
    \DisplayProof
\end{defish}

Bu temel düşünce $n$ değerli bir mantığa uygulandığında, her doğruluk değeri
için bir tane olmak üzere iki yerine $n$ konumlu bir sekant hesabı elde
edilir. $V=\{\False,\Undef,\True\}$ olan üç değerli bir mantıkta sekant
$\Gamma\mid\Pi\mid\Delta$ biçimindedir. Bu sekant !!a{valuation}
$\pAssign v$ altında ancak ve ancak bazı $!A\in\Gamma$ için
$\pValue{v}(!A)=\False$, bazı $!A\in\Delta$ için
$\pValue{v}(!A)=\True$ ya da bazı $!A\in\Pi$ için
$\pValue{v}(!A)=\Undef$ ise sağlanır. Dolayısıyla başlangıç sekantları
$!A\mid !A\mid !A$ her zaman sağlanır.

\end{document}
